\cotefile{70}
\scanpage{33}
% archive sheet number: 32
% top: matrix (16); the last two rows and lower central entries are overwritten.
\[
\sigma_i=
\left(\begin{array}{c c c c c c c c}
\overset{\unsure{i-1}}{\overline{\left(\begin{smallmatrix}1&&0\\&\ddots&\\0&&1\end{smallmatrix}\right)}}
 &0&u_{i1}&\unsure{v_{i1}|\nu_{i1}}&v_{i1}&\cdots&v_{i1}&0\\
 & &\vdots&\vdots&\vdots&&\vdots&\vert\\
 & &u_{i,i-1}&v_{i,i-1}&\unsure{v_{ii}}&\cdots&\unsure{v_{ii}}&0\\
 &0&u_{ii}\deleted{\unsure{1}}&v_{ii}\deleted{\unsure{+1}}&\unsure{v_{i,i+1}}&\cdots&v_{i,i+1}&0\\
0&1&u_{i,i+1}&v_{i,i+1}&v_{i,i+1}&&&\\
 &0&0&\deleted{\unsure{v_{i,i+2}}}\,0&\deleted{\unsure{v_{i,i+2}}}\,0&&0&0\\
 & & & &\deleted{\unsure{v_{i,i+2}}}\,0&&&\\
\unsure{\overline{0\ \cdots}}&\overline{0}&\overline{0}&\overline{0}&\overline{0}&&\overline{v_{i,i+2}}&\overline{0}\\
0\ \cdots\ \deleted{0}\,0& &0&0&0&&0&\unsure{1}
\end{array}\right)
\tag*{(16)}
\]

\[
\begin{aligned}
\sigma_i(x_1,\ldots,x_{n+1})={}&(x_1+u_{i1}x_{i+1}+v_{i1}\overset{\inserted{x_{i+2}+\cdots+x_{n+1}}}{\Sigma_i},\ x_2+u_{i2}x_{i+1}+v_{i2}\Sigma_i,\ \ldots\\
 &\cdots\ x_{i-1}+u_{i,i-1}x_{i+1}+v_{i,i-1}\Sigma_i;\\
 &u_{i,i}x_{i+1}+v_{ii}\Sigma_i,\ x_{i\deleted{\unsure{+1}}}+u_{i,i+1}x_{i+1}+v_{i,i+1}\Sigma_i;\\
 &v_{i,i+2}x_{i+2},\ \ldots,\ v_{i,i+2}x_{n+1})
\end{aligned}
\tag*{(17)}
\]

\[
H_i:\left\{\begin{aligned}
u_{i1}x_{i+1}+v_{i1}\Sigma_i&=u_{i2}\unsure{x_{i+2}|x_{i+1}}+v_{i2}\unsure{\Sigma_2|\Sigma_i}=\cdots=u_{i,i-1}x_{i+1}+v_{i,i-1}\Sigma_i=0\\
\inserted{-x_i+}u_{ii}x_{i+1}+v_{ii}\Sigma_i&=\deleted{\unsure{a}}0,\quad x_i+(u_{i,i+1}\deleted{\unsure{x_{i+1}}}-1)x_{i+1}+v_{i,i+1}\Sigma_i=0\\
(v_{i,i+2}-1)x_{i+2}&=\cdots\ (v_{i,i+2}-1)x_{n+1}=0
\end{aligned}\right.
\tag*{(18)}
\]

$\sigma_i$ est une \unsure{réflexion} ssi \unsure{ces relations sont conséquence} de la $i$ième (qui \unsure{exprimera} $x_i$ en \unsure{comb.} lin. de $x_{i+1},\ldots,x_{n+1}$. Ceci signifie:
\[
\left\{\begin{aligned}
u_{ii}+u_{i,i+1}&=1&v_{i,i}+v_{i,i+1}&=0\\
v_{i,i+2}&=1\\
u_{i1}=\cdots=u_{i,i-1}&=0,&v_{i1}&=\deleted{\unsure{0}}\cdots=v_{i,i-1}=0
\end{aligned}\right.
\tag*{(19)}
\]

donc les seuls paramètres qui \unsure{restent} \deleted{sont} $u_{i,i}\inserted{=\lambda_i}$ et $v_{i,i}=\mu_i$ \inserted{(ce dernier pour \unsure{$0\leq i\leq n-1$|$1\leq i\leq n-1$}, pas pour $\sigma_n$)}, de sorte que l'on peut écrire
\[
\begin{aligned}
\sigma_i(x_{\deleted{\unsure{0}}1},\ldots,x_{n+1})={}&(x_1,\ldots,x_{i-1},\ \lambda_i x_{i+1}+\mu_i\overset{\inserted{x_{i+2}+\cdots+x_{n+1}}}{\Sigma_i},\ x_i\deleted{\unsure{+}}+(1-\lambda_i)x_{i+1}-\mu_i\Sigma_i,\\
&x_{i+2},\ldots,x_{n+1})
\end{aligned}
\tag*{(20)}
\]
\scanpage{34}
% archive sheet number: 33
On a $(1\leq i\leq n)$
\[
\sigma_i^2(x_1,\ldots,x_{n+1})=(x_1,\ldots,x_{i-1};X_i,X_{i+1},x_{i+2},\ldots,x_{n+1})
\]
\[
X_i=\lambda_i\bigl[x_i+(1-\lambda_i)x_{i+1}-\mu_i\Sigma_i\bigr]+\mu_i\Sigma_i
=\lambda_i x_i+\lambda_i(1-\lambda_i)x_{i+1}
\]
\[
\begin{aligned}
X_{i+1}&=\bigl[\lambda_i x_{i+1}+\deleted{\mu_i\Sigma_i}\bigr]+(1-\lambda_i)\bigl[x_i+(1-\lambda_i)x_{i+1}-\mu_i\Sigma_i\bigr]-\deleted{\mu_i\Sigma_i}\\
&=(1-\lambda_i)x_i+x_{i+1}+(1-\lambda_i)\mu_i\Sigma_i
\end{aligned}
\]

On a
\[
\left.\begin{aligned}
X_i&\equiv x_i\quad\text{ssi}\quad\lambda_i=1\\
X_{i+1}&\equiv x_{i+1}\quad\text{ssi}\quad\lambda_i=1
\end{aligned}\right\}
\qquad\boxed{\sigma_i^2=\mathrm{id}\quad\text{ssi}\quad\lambda_i=1}
\tag*{(21)}
\]

\unsure{De|D°} $=$
\[
\sigma_0^2(x_1,\ldots,x_{n+1})=(X_1,x_2,\ldots,x_{n+1})
\]
\[
\begin{aligned}
X_1&=1+(\lambda_0-1)\bigl[1+(\lambda_0-1)x_1+(\mu_0-1)\Sigma_0\bigr]+(\mu_0-1)\Sigma_0\\
&=\lambda_0+(\lambda_0-1)^2x_1+\lambda_0(\mu_0-1)\Sigma_0
\end{aligned}
\]
\[
X_1\equiv x_1\quad\text{ssi}\quad\lambda_0=0\qquad\text{\unsure{i.e.|v.e.}}\quad
\boxed{\sigma_0^2=\mathrm{id}\quad\text{ssi}\quad\lambda_0=0}
\tag*{(22)}
\]

\[
\left\{\begin{aligned}
&\text{\uline{Cas général} paramètres}\quad
\left\{\begin{aligned}\lambda_i\deleted{\unsure{,\mu_i}}&&0\leq i\leq n\\\mu_i&&0\leq i\leq n-1\end{aligned}\right.\\
&\qquad\text{\uline{donc $2n+1$ paramètres}}\\[1ex]
&\text{\uline{Cas $n$-\unsure{polyèdral|polyédral}}}\quad(\sigma_i^2=1\quad\text{pour}\quad0\leq i\leq n)\quad
\begin{aligned}\lambda_0&=0\\\lambda_i&=1\quad1\leq i\leq n\end{aligned}\\
&\qquad\text{il reste $n$ paramètres}\quad\mu_i\ \ldots
\end{aligned}\right.
\]
\scanpage{35}
% archive sheet number: 34
\[
\left\{\begin{aligned}
\sigma_0(x,y)&=(1+(\lambda_0-1)x+(\mu_0-1)y,y)=(1+(\lambda_0-1)x+\alpha_0y,y)=(1+\unsure{\eta_0|\gamma_0}x+\alpha_0y,y)\\
\sigma_1(x,y)&=(\lambda_1y\deleted{\unsure{+}},x+(1-\lambda_1)y)
\end{aligned}\right.
\]
\margin{$k$ corps de car. $p>0$}

Partie homogène de $\sigma_0$:
\[
\begin{pmatrix}\deleted{\unsure{\lambda_0-1}}\eta_0&\alpha_0\\0&1\end{pmatrix}
\]
\deleted{\unsure{$\det(t\,\mathrm{id}-\sigma_{0V})=$}}

valeurs propres: $1$ et $\lambda_0-1=\eta_0$
\[
\deleted{\unsure{\sigma_{0V}}}\quad
\begin{pmatrix}a&c\\0&b\end{pmatrix}
\begin{pmatrix}a'&c'\\0&b'\end{pmatrix}
=\begin{pmatrix}aa'&ac'+b'c\\0&bb'\end{pmatrix}
\overset{\text{si }b'=1}{=}\begin{pmatrix}aa'&c+ac'\\0&b\end{pmatrix}
\]
\[
\begin{pmatrix}a&c\\0&b\end{pmatrix}
\begin{pmatrix}\eta_0&\alpha_0\\0&1\end{pmatrix}
=\begin{pmatrix}\eta_0a&c+\alpha_0a\\0&b\end{pmatrix}
\]
\[
\begin{pmatrix}\eta_0&\alpha_0\\0&1\end{pmatrix}^{2}
=\begin{pmatrix}\eta_0^2&\alpha_0+\deleted{\unsure{\alpha}}\unsure{\alpha_0\eta_0}\\0&1\end{pmatrix}
\]
\[
\begin{pmatrix}\eta_0&\alpha_0\\0&1\end{pmatrix}^{3}
=\begin{pmatrix}\eta_0^3&\alpha_0(1+\eta_0+\eta_0^2)\\0&1\end{pmatrix}
\]
\[
\overset{\unsure{-}}{\begin{pmatrix}\eta_0&\alpha_0\\0&1\end{pmatrix}}^{\unsure{n|m}}
=\begin{pmatrix}\eta_0^n&\alpha_0(1+\eta_0+\cdots+\unsure{\eta_n^{n-1}|\eta_0^{n-1}})\\0&1\end{pmatrix}
\]
\[
\sigma_{0V}^n=1\quad\Longleftrightarrow\quad
\left\{\begin{aligned}
\eta_0^n&=1\quad\text{i.e.}\quad(\unsure{\eta_0|\eta}-1)(1+\eta_0+\cdots+\eta_0^{n-1})=0\\
\deleted{\unsure{\alpha_0(1+}}&&\alpha_0(1+\eta_0+\cdots+\eta_0^{n-1})=0
\end{aligned}\right.
\]
\[
\Updownarrow\qquad\sigma_0^n\text{ une translation}
\]

Deux cas
\[
\left\{\begin{aligned}
\text{a)}\quad&\deleted{\unsure{\eta_0^n}}1+\eta_0+\cdots+\eta_0^{n-1}=0\quad\text{i.e.}\quad
\deleted{\unsure{(\eta_0\ne1)|(\eta_0^n=1)}}\eta_0^n=1\\
&\hspace{7em}\text{et }(\eta_0\ne1\text{ ou }p\deleted{\unsure{\nmid}}\mid n)\\
\text{b)}\quad&\eta_0=1,\ \alpha_0=0\quad[\text{i.e. }\sigma_{0V}=\mathrm{id}\text{ i.e. }\sigma_0\text{ est la}\\
&\qquad\text{translation par }(1,0)]\text{ et }p\nmid n
\end{aligned}\right.
\]
\[
\sigma_0^n(0)=?\qquad
\begin{aligned}
&\deleted{\sigma_0(0)=(1,0)}\\
&\deleted{\unsure{\sigma_0^2(0)=(1+\eta_0}}
\end{aligned}
\]
\[
\begin{aligned}
\sigma_0^{\deleted{\unsure{1}}}(x,0)&=(1+\eta_0x,0)\\
\sigma_0^2(x,0)&=1+\eta_0+\eta_0^2x\\
\sigma_0^3(x,0)&=1+\eta_0+\eta_0^2+\eta_0^3x\\
&\cdots\\
\sigma_0^n(x,0)&=1+\eta_0+\cdots+\eta_0^{n-1}+\eta_0^n x
=\left\{\begin{aligned}
&\text{\unsure{Cas a)}}\quad\eta_0^n x=x\\
&\text{\unsure{Cas b)} translation par }(n,0)\\
&\qquad\text{qui n'est pas l'identité.}
\end{aligned}\right.
\end{aligned}
\]
\scanpage{36}
% archive sheet number: 35
On trouve donc
\[
\sigma_0^{\unsure{\ell|\varepsilon}}=\mathrm{id}\quad\Longleftrightarrow\quad
1+\unsure{\eta_0|\gamma_0}+\cdots+\eta_0^{\ell-1}=0
\]
donc

Ordre de $\sigma_0=$ ordre de $\eta_0\in k^*$ \uline{sauf} si $\eta_0=1$ \inserted{i.e. $\det\sigma_{0V}=1$}, auquel cas $\sigma_{0V}$ est la transvection
$\begin{pmatrix}1&\alpha_0\\0&1\end{pmatrix}$
(qui est d'ordre $1$ ssi $\alpha_0=0$, sinon d'ordre $p$)
\[
\eta_0=1\quad\left\{\begin{aligned}
\text{a)}\quad&\alpha_0=1\quad\sigma_0=\mathrm{id}\qquad\text{\uline{cas impropre}}\\
\text{b)}\quad&\alpha_0\ne1\quad\left\{\begin{aligned}
p=0\quad&\sigma_{0V}\text{ d'ordre infini, a fortiori }\sigma_0\text{ d'o. inf.}\\
p>0\quad&\sigma_{0V}\text{ d'ordre }p,\quad\sigma_0^p=\deleted{\text{translation}}\inserted{\text{\unsure{au pt.}}}\\
&\deleted{\unsure{(n,0)}}\quad\left\{\begin{aligned}
&\deleted{p\mid n\ \Longrightarrow}\ \sigma_0^p=\mathrm{id}\\
&\deleted{p\nmid n\quad\sigma_0\text{ d'ordre }p^2}
\end{aligned}\right.
\end{aligned}\right.
\end{aligned}\right.
\]
\par\noindent\rule{.4\linewidth}{.4pt}\par
\[
\begin{aligned}
\sigma_1^2(x,y)&=(\lambda_1x+\lambda_1(1-\lambda_1)y,(1-\lambda_1)x+(\lambda_1+(1-\lambda_1)^2)y)\\
&=(\lambda_1x+\lambda_1(1-\lambda_1)y,(1-\lambda_1)x+(\lambda_1^2-\lambda_1+1)y)
\end{aligned}
\]
\[
\sigma_1=\begin{pmatrix}0&\lambda_1\\1&1-\lambda_1\end{pmatrix}
\]
\[
\det(t\,\mathrm{id}-\sigma_1)=t^2-(1-\lambda_1)t-\lambda_1
=\deleted{t^2-(1+\delta)t+\delta}
\]

Valeurs propres \deleted{\unsure{$\eta,\eta'$}}
\inserted{$1,-\lambda_1=\deleted{\delta}\eta_1$}
\[
\deleted{\left\{\begin{aligned}-\lambda_1&=\eta\eta'\\1-\lambda_1&=\eta+\eta'\end{aligned}\right.}
\qquad\deleted{-\lambda_1=\delta_1=\delta}
\]
\[
\deleted{\left\{\begin{aligned}\eta\eta'&=\delta\\\eta+\eta'&=1+\delta\end{aligned}\right.}
\]
% margin note written slanted
\margin{\unsure{alors|choix}}
\[
\deleted{\eta^{\unsure{n|m}}=\eta'{}^n=1}\quad\deleted{\Longrightarrow\delta^n=1}
\]
\[
\deleted{\frac{\eta}{\delta}}\qquad\deleted{\frac{\eta'}{\delta}}\qquad
\deleted{\frac{\eta}{\delta}}\qquad\deleted{\frac{\eta'}{\delta}}\qquad
\deleted{y=}\quad x=\lambda_1y
\]
\[
\deleted{\unsure{\delta\ne1}}
\]

Ordre de $\sigma_1=$ ordre de $\eta_1\in k^*$, \uline{sauf si $\delta=1$} i.e. $\lambda_1=-1$
\[
\sigma_1(x,y)=(-y,x+2y)\quad\text{transvection}
\]
\unsure{\deleted{non}|$\ne$} identité, d'ordre infini si $p=0$, d'ordre $p$ sinon.
\scanpage{37}
% archive sheet number: 36
\[
\begin{array}{|l}
\sigma_0(x_1,\ldots,x_{n+1})=(1+\unsure{\eta_0|\gamma_0}x_1+\alpha_0\Sigma_0,x_2,\ldots,x_{n+1})\\
\quad\cdots\\
\begin{aligned}
1\leq i\leq n-1\qquad\sigma_i(x_1,\ldots,x_{n+1})={}&(x_1,\ldots,x_{i-1},-\eta_i x_{i+1}+(1+\alpha_i)\Sigma_i,\ x_i+(\eta_i+1)x_{i+1}-(1+\alpha_i)\Sigma_i,\\
&x_{i+2},\ldots,x_{n+1})
\end{aligned}\\
\quad\cdots\\
\sigma_{n\deleted{\unsure{+1}}}(x_1,\ldots,x_{n+1})=(x_1,\ldots,x_{n-1},-\eta_nx_{n+1},x_n+(1+\eta_n)x_{n+1})\\\hline
\end{array}
\]

Valeurs propres de $\sigma_i$:
\[
\left\{\begin{aligned}&1\text{ de multiplicité }n-1\\&\eta_i\end{aligned}\right.
\]

ordre de $\sigma_i\overset{?}{=}$ ordre de $\eta_i\in k^*$ si $\eta_i\ne1$
\scanpage{39}
% archive sheet number: 38
\margin{\unsure{(9 p)}}

\deleted{\unsure{Th. des Topos|Th. de Topos}}\qquad\deleted{\unsure{7}}

``étude locale'' des $n$-\unsure{polyèdres} et des $n$-cartes

(vieille rédaction)
\scanpage{40}
% archive sheet number: 39
Soit \deleted{\unsure{$X$}} une carte combinatoire \inserted{$C$} de moquette $E$ --- un $G_n$-ensemble --- où
\[
\begin{aligned}
G_n=\{\sigma_0,\ldots,\sigma_n\mid{}&\sigma_0^2=\cdots=\sigma_n^2=1,\\
&\sigma_i\sigma_j=\sigma_j\sigma_i\quad\text{si }j\geq i+2\}
\end{aligned}
\]

Rappelons que les $i$-facettes de la carte sont les éléments de l'ens
\[
D_i=E/(\sigma_0,\ldots,\widehat{\sigma_i},\ldots,\sigma_n)
=E/(G_n^{\deleted{\unsure{+}}-}(i)\times G_n^+(i))
\]
\inserted{$G_n^-(i)\subset G_n$, $G_n^+(i)\subset G_n$ sont engendrés respectiv. par $\sigma_0,\ldots,\sigma_{i-1}$ et $\sigma_{i+1},\ldots,\sigma_n$.}

Considérons, pour $f_i\in D_i$ fixé, l'ens des facettes \inserted{$f'$} de $C$ \deleted{\unsure{dont un qui}} $<f_i$, \deleted{donc} i.e. de dim $j<i$ et telles que $(f'_j,f_i)$ \unsure{provienne|proviennent} d'un él. de $E$. Soit \deleted{\unsure{$E(i)^+$}}
\[
E^{\deleted{+}-}(i)=E/G_n^{\deleted{-}+}(i)\longrightarrow D_i
\]
\deleted{\unsure{``$D$ augmenté i''}}
\deleted{\unsure{$E^+(i)=E/G_n$}}

c'est un ens. à opér. $G_n^+(i)\simeq G_{i-1}$,
\deleted{et la \unsure{flèche vers|flèche sur} $D_i$ est \unsure{compatible} \unsure{avec ces|avec les} \unsure{opérations}}
\deleted{\unsure{$D_i$ est un quotient.}}
et sauf erreur c'est un \unsure{$(i-1)$-un. polyèdre} \inserted{\unsure{pur;|faux;}} i.e. $\sigma_0,\ldots,\sigma_{i-\unsure{1}}$ et les produits de tels éléments qui commutent opèrent sans pts fixes dans $E^-(i)$.
\uline{\unsure{Ainsi} $E^-(i)$ \unsure{les fibres de}} \inserted{\uline{correspond:}} \deleted{\unsure{est}} une $(i-1)$-carte combinatoire, dont les comp. connexes correspondent aux él. de $D_i$ i.e. aux $i$-facettes $f_i$.
Soit $\Pi_{f_i}^{\unsure{-}}$ \inserted{la \deleted{\unsure{poly}} $(i-1)$-carte connexe \deleted{qui} \unsure{déf.} par $f_i$} \deleted{la fibre de $f_i$}. Ces $j$-facettes sont
\scanpage{41}
% archive sheet number: 40
les él. de
\[
E^-_{\deleted{\unsure{f}}}(i)_{f_i}/(\sigma_0,\ldots,\widehat{\sigma_j},\ldots,\sigma_{i-1}),\quad\text{i.e.}
\]
les él. de
\[
E^-(i)/(\sigma_0,\ldots,\widehat{\sigma_j},\ldots,\sigma_{i-1})
\simeq E/(\sigma_0,\ldots,\widehat{\sigma_j},\ldots,\widehat{\sigma_i},\ldots,\sigma_n)
\]
qui \deleted{\unsure{au delà}} s'envoient dans $f_i\inD_i$, i.e. les drapeaux de types $(j,i)$ dont la composante de \deleted{degré} dim $i$ est $f_i$. Plus gén, ses drapeaux de type $(j_1<\cdots<j_p)$ (i.e. $0\leq j_1,\deleted{\unsure{j}}j_p\leq\unsure{i-1|i}$) sont les drapeaux de type $(j_1,\ldots,j_p,i)$ de $\Pi$ \deleted{dont} la composante de degré $i$ est $f_i$. On l'appelle \unsure{le|la} \inserted{\uline{carte}} ``\uline{bord}'' de la facette $f_i$, \unsure{notée|noté} $\partial f_i$.

\unsure{Dualement|Duellement}, considérons
\[
E^+(i)=E/G_n^-(i)\inserted{=D_{i,i+1,\ldots,n}}\longrightarrow D_i
\]
sur lequel $G_n^+(i)\simeq\deleted{\unsure{n-i}}G_{n-i-1}$ opère, le quotient étant $D_i$. On trouve une moquette \inserted{de $(n-i-1)$-cartes} (sauf erreur), dont les facettes de dim $j$ $(0\leq j\leq n-i-1)$ ``sont'' les drapeaux de type $(i,i+j+1)$ \inserted{de $\Pi$} \deleted{dont} la composante de dim $i$ est $f_i$.

On peut aussi, pour \deleted{deux facettes} \inserted{un drapeau \inserted{$d_{ij}$} de type $(i,j)$, \unsure{défini par|définissant}}
\[
f_i<f'_j\qquad(\deleted{i<j}\inserted{j\geq i+2})
\]
considérer
\[
\deleted{\unsure{E(i,j)=E/(\sigma_0-\sigma_{i-1};\sigma_{j+1},\ldots,\sigma)}}
\]
\scanpage{42}
% archive sheet number: 41
\[
E(i,j)=E/(\sigma_0,\ldots,\sigma_{i-1};\sigma_{j+1},\ldots,\sigma_n)
=E/G^-(i)\times G^+(j)=D_{i,i+1,\ldots,j}
\]
et l'application
\[
E(i,j)\longrightarrow D_{i,j}\ \deleted{\unsure{D_j}}
\]
\margin{$0\leq i\leq i+2\leq j\leq\unsure{n+1}$}

Le groupe engendré par $\sigma_{i+1},\ldots,\sigma_{j-1}$ opère sur $E(i,j)$, avec comme quotient $D_{i,j}$. Sans doute il opère \uline{cartographiquement}, et on trouve une $(j-i-2)$-carte, dont les composantes connexes correspondent aux $d_{ij}\in D_{ij}$; \inserted{sur $\Pi(d_{ij})$} les facettes de dim $h$ $(0\leq h\leq j-i-2)$ correspondent aux drapeaux de type $(i,\underbrace{i+h+1}_{\inserted{i+1\leq\ \leq j-1}},j)$ au dessus du drapeau $d_{ij}$.

Cas particuliers:
\[
\begin{aligned}
j=i+2&\quad\text{\unsure{un peu idiot}}\\
j=i+3&:\quad\text{on trouve un polygône, d'où}\\
&\qquad\text{un \uline{invariant}}\quad\unsure{p|\rho}(d_{ij})\in\overline{\mathbb{N}}^*
\end{aligned}
\]
\[
(\text{où }1\leq p(d_{ij})\leq+\infty)
\]

\deleted{On \unsure{s'intéresse au}}
\deleted{Introduisons}
Formellement les $\Pi^{\unsure{-|=}}(f_i)$ correspondent à \deleted{\unsure{au}} $\Pi(d_{-1,i\deleted{\unsure{f_i}}})$, $d_{-1,i}$ défini par $f_i$ \deleted{\unsure{pour $\Pi$}} \deleted{\unsure{$(-1)$-facettes}} et $\Pi^+(\unsure{f^i|f_i|f'})$ \deleted{par} à $\Pi(d_{i,n+1})$, où $d_{i,n+1}$ est défini par $f_i$. Les invariants numériques
\scanpage{43}
% archive sheet number: 42
essentiels sont

a) Les $\unsure{\nu|v}(\Pi^-_{f_2})$ --- \unsure{dénoté} $\unsure{p|\rho}_0(\Pi)$ si $\Pi$ \uline{régulier}

b) les $\nu(\Pi(d_{0,3})),\nu(\Pi(d_{1,4}))\cdots(\nu(\Pi(d_{n-3,n})))$

notés $p_1(\Pi),p_2(\Pi),\ldots,p_{n-2}(\Pi)$ si $\Pi$ \uline{régulier}

c) les $\nu(\Pi^+_{f_{n-2}})$, \unsure{dénoté|noté} $p_{n-1}(\Pi)$ si $\Pi$ \uline{régulier}

\par\noindent\hspace*{.25\linewidth}\rule{.25\linewidth}{.4pt}\par

Supposons $\Pi$ \uline{régulier}, je dis que
\[
\left\{\begin{aligned}
p_0&=\text{ordre de }\sigma_0\sigma_1\\
p_1&:\text{ordre de }\sigma_1\sigma_2\\
&\cdots\\
p_{n-1}&:\text{Ordre de }\sigma_{n-1}\sigma_n
\end{aligned}\right.
\]
\margin{vrai tout au moins dans le cas des géométries d'incidences (``polyèdres''\ldots)}

En effet, \deleted{\unsure{$\Pi$}} sous hyp. de \uline{régularité}, $\nu(\Pi^-_{f_2})$ est l'ordre de $\sigma_0\sigma_1$ opérant sur $E^-(2)_{f_2}\inserted{=(E/G_n^+(2))_{f_2}}$, \deleted{\unsure{$\nu(\Pi(d^{\unsure{a|o}}_{i,i+3}))$ est l'ordre}}
\[
(\sigma_0\sigma_1)^p x=\deleted{\unsure{x}}g_x\deleted{\unsure{\cdot}}x\qquad g_x\in G_n^+(2)
\]
\diagram{\(\begin{array}{ccccc}
i-1&&\unsure{i+1\leq j-1}&j+1&\\
&i&j&&i\\
&&\unsure{\geq|=}&&\\
&&i+2&&
\end{array}\)}
\scanpage{44}
% archive sheet number: 43
% upper right corner of the sheet is torn.
$\Pi$ une $n$-carte \inserted{comb.} de moquette $R=\operatorname{Rep}(\Pi)$, $G_n$-ensemble (où
\[
G_n=\{\sigma_0,\ldots,\sigma_n\mid\sigma_0^2=\cdots=\sigma_n^2=1,\ \sigma_i\sigma_j=\sigma_j\sigma_i\text{ si }j\geq i+\unsure{2}\}
\]

Une \uline{représentation géom} $\unsure{\varphi|\phi}$ de $\Pi$ dans \deleted{\unsure{de rg $n+1$}} un espace affine $E$ (sur \unsure{anneau} $k$ disons) est la donnée d'\deleted{\unsure{une hom.}}application
\diagram{(1)\quad\begin{tikzcd}[ampersand replacement=\&]
R\arrow[r,"\varphi_R"]\&\operatorname{Drap\,max}(E)\arrow[d,equal]\\
\&\{(E_0\subset E_1\cdots\subset E_i\subset\cdots\subset E_n\subset E)\}
\end{tikzcd}}
\inserted{$E_i$: ss-espace affine de $E$ de rg $i$}

telle que pour $r\in R$, $0\leq i\leq n$, on ait
\[
\left\{\begin{aligned}
\pi_j(\varphi_R(\sigma_i r))&=\pi_j\varphi_R(r)\quad\text{si}\\
0\leq i,j&\leq n,\qquad i\ne j
\end{aligned}\right.
\tag*{(2)}
\]
où $\pi_j(d)$ désigne la composante de dim $j$ d'un drapeau $d$. Il revient au \unsure{$=$|même} \inserted{\unsure{de par.|de dire}} à dire que par passage au quotient, $\pi_j\varphi_R$ se factorise en
\diagram{(3)\quad\begin{tikzcd}[ampersand replacement=\&]
D_j(\Pi)\arrow[r,"\varphi_j"]\&\unsure{\mathrm{Gr\,aff}_j|\mathrm{Grass\,aff}_j}(E)
\end{tikzcd}\quad$(0\leq j\leq n)$}
\scanpage{45}
% archive sheet number: 44
% Sheet read after rotation through 180 degrees; the entire draft is crossed out.
\margin{\deleted{(\unsure{régulière équivar.})}}
\deleted{Considérons une réalisation géom. $\unsure{\varphi|\phi}$ de $\Pi$ \inserted{dans l'espace affine $E$ de rg $n+1$} (\unsure{affine dans}) sur un anneau $k$, donnée par}

\diagram{\begin{tikzcd}[ampersand replacement=\&,row sep=large,column sep=huge]
\deleted{G=\operatorname{Aut}(\Pi)} \arrow[r,"\deleted{\varphi_G}"] \& \deleted{\operatorname{Aff}(E)} \\
\deleted{R=\operatorname{Rep}(\Pi)} \arrow[r,"\deleted{\varphi_R}"] \arrow[d] \arrow[dr,phantom,"\deleted{\text{$G$-hom.}}"] \& \deleted{\operatorname{Drap\ max}(E)} \arrow[d,"\deleted{\pi_i}"] \\
\deleted{D_i(\Pi)} \& \deleted{\unsure{\operatorname{Gr.aff}|\operatorname{Gr.\ aff}}_i(E)}
\arrow[from=3-1,to=3-2,"\deleted{\varphi_{D_i}}"]
\arrow[from=2-1,to=3-1,phantom,"\deleted{\text{\unsure{$G$-hom.|$G$-homo}}}"',shift left=6ex]
\end{tikzcd}}

\[
\left\{\begin{aligned}
&\deleted{\varphi_R(r)=(E_0\subset E_1\subset\cdots\subset E_n\subset E)}\\
&\deleted{\sigma_i E_j=E_j\quad\text{si }i\ne j\qquad\ldots}
\end{aligned}\right.
\]
\scanpage{46}
% archive sheet number: 45
Lorsque la $n$-carte $\Pi$ est un polyèdre (= géométrie d'incidences), alors la donnée de $\unsure{\varphi|\phi}_R$ satisfaisant (2) \uline{équivaut} à la donnée des $\varphi_j$ $(0\leq j\leq n)$, \unsure{soumises|soumis} à la condition d'être compatibles avec les relations d'incidences. Dans le cas général, on peut dire \inserted{encore} que (3) se généralise en des applications

\diagram{\begin{tikzcd}[ampersand replacement=\&,column sep=huge]
(4)\quad \unsure{D|\mathscr{D}}_{\{j_0,j_1,\ldots,j_p\}}(\Pi) \arrow[r,"\varphi_{\{j_0,\ldots,j_p\}}"] \& \operatorname{Drap\ aff}_{\{j_0,j_1,\ldots,j_p\}}(\Pi)
\end{tikzcd}}

pour \unsure{tte} partie $\{j_0<j_1<\cdots<j_p\}$ de $\Delta_n:\{0,1,\ldots,n\}$, ces applications étant compatibles avec les applications verticales pour des inclusions de parties de $\Delta_n$. La donnée de $\varphi_R$ satisfaisant (2) équivaut à la donnée des applications (4) satisfaisant la condition de fonctorialité précédente --- on peut donc les interpréter \inserted{les $\varphi_R$} comme des hom. de géométries de drapeaux basées sur \inserted{l'ens\unsure{ble} ordonné} $\Delta_n$, de $D_{\unsure{*|\bullet}}(\Pi)$ dans $\operatorname{Drap.aff}_*(E)$.
\scanpage{47}
% archive sheet number: 46
% The first draft, through ``i.e. un'', is cancelled by long diagonal strokes.
\deleted{Ceci dit, considérons la $r=(j-i-2)$-carte \inserted{$\Pi_{i,j}$} définie par la moquette $(0\leq i\leq i+2\leq j\leq n)$}
\[
\deleted{R_{ij}=R/(\sigma_0\cdots\sigma_{i-1},\sigma_{j+1},\ldots,\sigma_n)}
\]
\deleted{(\uline{NB}}
\[
\deleted{R_{0j}=R/(\sigma_{j+1},\ldots,\sigma_n),\qquad R_{i,n}=R/(\sigma_0\cdots\sigma_{i-1}))}
\]
\deleted{avec les opérations induites par passage au quotient par}
\diagram{$\begin{array}{ccc}
\deleted{\sigma_{i+1}}&\deleted{,\ldots,}&\deleted{\sigma_{j-1}}\\
\deleted{\Vert}&&\deleted{\Vert}\\
\deleted{\sigma'_0}&\deleted{,\ldots,}&\deleted{\sigma'_{j-i-2=r}}
\end{array}$}
\deleted{qui en font une $r$-carte, dont l'ens des composantes connexes s'identifie à $D_j(\Pi)$. Fixons une telle composante connexe, i.e. un}

Ceci dit, soit $f_i\in D_i(\Pi)$ \inserted{$(1\leq i\leq n)$}, $E_i=\unsure{\varphi|\phi}_i(f_i)$, et considérons la $(i-1)$-carte induite $\unsure{\partial|\delta}f_i^{\unsure{\subset|c}}$ dont l'ens des repères est $D_{01\ldots i}(f_i)$ = image rec. de $\{f_i\}$ par
\diagram{\begin{tikzcd}[ampersand replacement=\&]
D_{0,1,\ldots,i}\inserted{(\Pi)} \arrow[r,phantom,"{\unsure{-|\longleftarrow|\longrightarrow}}" description] \& D_i(\Pi).
\end{tikzcd}}
On trouve alors \inserted{en termes de $\varphi$} une réalisation géom. \unsure{canoniq.} (dite ``induite'') de $\partial f_i$ dans $E_i$.

Légèrement moins évidente est la définition de la représentation induite de la $(n-i-1)$-carte $\unsure{\partial f_i^{\supset}|\partial f_i\supseteq}$ dont l'ens des
\scanpage{48}
% archive sheet number: 47
repères est $D_{i,i+1,\ldots,n}(f_i)=\text{image inv. de }\{f_i\}$ par
\diagram{\begin{tikzcd}[ampersand replacement=\&]
D_{i,i+1,\ldots,n} \arrow[r] \& D_i.
\end{tikzcd}}
\deleted{\unsure{Ce}} Les sommets sont les \deleted{$\unsure{f_{i+1}|f_{1n}}$} $(i,i+1)$-drapeaux de $\Pi$ dont la \unsure{$i$-composante|$i$\textsuperscript{ème} composante} est $f_i$, et ils s'envoient sur des drapeaux de $E$ de type $(i,i+1)$ dont la $i$-composante est $f_i$, \deleted{\unsure{les}} i.e. des couples $(E_{i\deleted{\unsure{2}}}\subset E_{i+1})$, $E_{i+1}$ de rg $i+1$.

Mais si $V_i$ est l'espace des translations de $E_i$, les $E_{i+1}$ correspondent à des droites de $\overline{E}=E/V_i$.
